
Selecting the LoRA rank for each transformer layer is typically treated as a hyperparameter problem, requiring grid search, gradient signals, or calibration runs before a single fine-tuning step begins. Yet the pretrained weight matrices that will be adapted already encode, in their singular value spectra, a structural fingerprint of how much representational complexity each layer developed during pretraining. If that fingerprint predicts how much rank a layer needs to absorb task-relevant signal, rank selection could be resolved in a single SVD pass --- before any training data is observed.

Low-Rank Adaptation (LoRA) \citep{Hu2022LoRA} achieves parameter-efficient fine-tuning by decomposing weight updates as $\Delta W = BA$ where $B \in \mathbb{R}^{d \times r}, A \in \mathbb{R}^{r \times k}$, and rank $r \ll \min(d, k)$. Its central weakness is the uniform rank assumption: every layer receives the same rank $r$, regardless of whether that layer's contribution to fine-tuning is large or small. This assumption is known to be suboptimal: Aghajanyan et al. [2021] showed that intrinsic dimensionality varies substantially across transformer layers, and LAARA \citep{Tripathi2026LAARA} formally proves that uniform rank allocation is inferior to any Fisher-optimal per-layer assignment.

The deeper problem is that all existing per-layer rank methods require some form of training or calibration signal. Gradient-based methods --- AdaLoRA \citep{Zhang2023AdaLoRA}, DyLoRA \citep{Valipour2022}, LAARA \citep{Tripathi2026LAARA} --- modify rank during training using importance scores derived from weight updates. Calibration-based methods --- IFCLoRA \citep{Zhang2026IFCLoRA} --- run a forward pass over a calibration dataset before fine-tuning. Both families incur overhead that is absent if rank can be predicted from the pretrained weights alone.

\textbf{Our central question:} Does the effective rank of a pretrained weight matrix --- $\text{erank}(W_0) = \exp(H(\sigma/\|\sigma\|_1))$, where $H$ is Shannon entropy and $\sigma$ are singular values --- correlate significantly with the rank assigned by an exhaustive per-layer oracle sweep?

We introduce the PARA (Per-layer Argmax Rank Assignment) oracle protocol and report correlation results for BERT-base-uncased. The experimental scope is one of three planned model families; multi-family results are pending.

Our contributions are:

\begin{enumerate}
\item \textbf{Empirical correlation test (H-E1):} Pearson $r = 0.984$ ($p = 0.0013, n = 5$) between erank($W_0$) and PARA oracle ranks for BERT-base-uncased, substantially exceeding the pre-registered threshold of $r \geq 0.65$.
\end{enumerate}

\begin{enumerate}
\item \textbf{Participation ratio agreement (P5):} Pearson $\rho = 0.968$ between erank and participation ratio PR($W_0$), confirming metric robustness.
\end{enumerate}

\begin{enumerate}
\item \textbf{erank-oracle correlation protocol:} Open implementation of erank computation, PARA oracle extraction, and correlation analysis, enabling replication and extension to additional model families.
\end{enumerate}

The paper proceeds as follows: Section 2 surveys related work. Section 3 describes the methodology. Section 4 details the experimental setup. Section 5 presents results for BERT-base-uncased. Section 6 discusses implications, limitations, and open questions. Section 7 concludes.

